\chapter{Choose a project and something to change}
\label{ch:choose}

\section{Start with something you already use}

The best first project is one you use, because you know where it confused you. If it is a tool for your field, you also understand its users. Look at its page on GitHub.

\begin{figure}[H]\centering\includegraphics[width=\linewidth]{web_oss_repo.png}\caption{A project's page on GitHub. The tabs at the top lead to its code, its issues and its pull requests. On the right are its description and licence. The \textbf{Fork} button sits above the file list.}\end{figure}

\section{Check that the project welcomes changes}

Before you spend time, check four things on that page. The date beside the most recent files shows whether anyone has been active this year. A \textbf{Licence} on the right shows others are allowed to reuse the code, and open-source projects must have one. A file named \texttt{CONTRIBUTING.md} sets out the project's rules, so read it first if there is one. And the \textbf{Pull requests} tab shows whether outsiders' changes get answered.

\section{Find something small}

Press the \textbf{Issues} tab, which lists known problems and requests. Many projects label easy ones \texttt{good first issue}.

\begin{figure}[H]\centering\includegraphics[width=\linewidth]{web_oss_issues.png}\caption{The Issues tab. This project has none open, which is common for small ones. In that case, improve something you noticed yourself.}\end{figure}

If you plan a large change, first write an issue describing it and wait for a reply, so that you do not spend days on something the maintainers would not accept. For a spelling fix, you can go straight ahead. If you take an existing issue, leave a comment saying you are working on it, so that nobody duplicates the work.

A useful first change is one a reader could check in a minute. Fix a broken link, correct a mistaken command in the instructions, or add a missing step you had to work out yourself.
